\scanpage{12}
\uline{Remarques} 1.11. a)Le préschéma $\handwritten{\mathscr T}$ de 6.10. s'appelera comme de juste le \uline{préschéma des tores maximaux} de $G$. On verra au N\textsuperscript{o}5 qu'il est en fait \uline{affine} sur $S$. On peut le vérifier directement, lorsque $G$ est isomorphe à un sous-préschéma en groupes fermé d'un préschéma de la forme $\mathrm{Gl}(n)$, en utilisant Exp XI 4.6. et en remarquant que par le raisonnement de la démonstration de 1.7. b) , $\handwritten{\mathscr T}$ s'identifie à un sous-préschéma à la fois \uline{ouvert et fermé} du préschéma $F$ qui représente les sous-groupes de type multiplicatif de $G$.

% A proofreader's slash crosses the space between des and tores in the underlined expression above.

b) Il \deleted{ne serait pas difficile} \inserted{\handwritten{est possible}} de donner une démonstration de 1.7. donc de 1.9. n'utilisant pas les résultats délicats de Exp XI, mais seulement \inserted{l\handwritten{es} \deleted{\unsure{faciles}} \handwritten{r}ésultats faci\handwritten{le}s \deleted{\unsure{de Exp XI N\textsuperscript{o}6}}} Exp XI 3.12. \inserted{\handwritten{et 6.2.,}} en travaillant uniquement avec des groupes de type multiplicatif finis sur la base (moralement, les groupes ${}_nT$ où $T$ est un tore maximal) , comparer \deleted{Exp} N\textsuperscript{o}7 . \deleted{Par contre,} \inserted{\handwritten{La même remarque vaut pour}} la démonstration de 1.10.\handwritten{,}

\margin{\handwritten{et 6.2.,}}

\deleted{s'appuie de façon essentielle sur Exp XI 4.1. Dans un autre exposé, nous donnerons une version généralisée de 1.10. par une démonstration différente (s'appuyant sur un théorème de représentabilité de EGA \unsure{V)|VI)}.}

\deleted{\deleted{Exception faite du \unsure{§, N\textsuperscript{o}5|N\textsuperscript{o}5}}, \inserted{\handwritten{De façon générale,}} \inserted{pour les résultats du N\textsuperscript{o}5,} les théorèmes de représentabilité de Exp XI N\textsuperscript{o} 4 \deleted{\unsure{me semblent}} pourrai\deleted{\unsure{ent}}\handwritten{t}\deleted{\unsure{on}} être éliminés de la théorie présentée \inserted{\handwritten{dans le présent exposé et les suivants}}}

\deleted{\handwritten{de façon générale, il est possible (et sans doute avantageux) d'éliminer totalement les résultats de représentabilité de Exp $\overline{\underline{\mathrm{XI}}}$ N\textsuperscript{o} 4 de la théorie présentée dans le présent exposé et les suivants.}}

% archive sheet number: 11
